% Journal-agnostic article scaffold.
% Compile with:  tectonic main.tex
% or classic:    pdflatex main && bibtex main && pdflatex main && pdflatex main
\documentclass[11pt]{article}

% --- Core packages (all ship with standard TeX Live / tectonic) ---
\usepackage[utf8]{inputenc}
\usepackage[T1]{fontenc}
\usepackage{amsmath}      % equations: equation, align, \eqref
\usepackage{amssymb}      % extra math symbols
\usepackage{graphicx}     % \includegraphics
\usepackage{booktabs}     % \toprule \midrule \bottomrule
\usepackage[numbers]{natbib} % \citep \citet ; numbered style
\usepackage[margin=1in]{geometry}
\usepackage{hyperref}     % clickable refs/links (load near last)
\hypersetup{colorlinks=true, linkcolor=black, citecolor=black, urlcolor=blue}

% Tell graphicx where images live (create a figures/ folder).
\graphicspath{{figures/}}

% --- Front matter ---
\title{A Concise, Descriptive Paper Title}
\author{%
  First Author\thanks{Corresponding author: \texttt{first.author@example.org}}\\
  \small Department, Institution\\
  \and
  Second Author\\
  \small Department, Institution%
}
\date{\today}

\begin{document}
\maketitle

\begin{abstract}
State the problem, what you did, and what you found in one paragraph.
Keep it to roughly 150--250 words (check your venue's limit).
Avoid citations and undefined acronyms in the abstract.
\end{abstract}

\noindent\textbf{Keywords:} keyword one; keyword two; keyword three

% =====================================================================
\section{Introduction}
\label{sec:intro}
Introduce the context and motivation, state the gap in prior work, and list
your contributions. Cite related work parenthetically \citep{smith2020example}
or in the text as \citet{jones2019book}. Close with a roadmap of the paper:
Section~\ref{sec:methods} describes the methods, Section~\ref{sec:results}
reports results, and Section~\ref{sec:discussion} discusses them.

% =====================================================================
\section{Methods}
\label{sec:methods}
Describe your approach in enough detail to reproduce it. Define notation here.
A numbered equation looks like this:
\begin{equation}
  \label{eq:loss}
  \mathcal{L}(\theta) = \frac{1}{N} \sum_{i=1}^{N} \ell\big(f_\theta(x_i),\, y_i\big).
\end{equation}
Refer to it later as Equation~\eqref{eq:loss}. Use inline math like
$\theta \in \mathbb{R}^d$ for short expressions.

% =====================================================================
\section{Results}
\label{sec:results}
Report findings without interpretation. Reference the table and figure below.

Table~\ref{tab:results} summarizes the main quantitative results.

\begin{table}[t]
  \centering
  \caption{Main results. Best value per column in bold.}
  \label{tab:results}
  \begin{tabular}{lcc}
    \toprule
    Method & Accuracy ($\%$) & Time (s) \\
    \midrule
    Baseline      & 81.2 & 12.4 \\
    Ours          & \textbf{88.7} & \textbf{9.1} \\
    \bottomrule
  \end{tabular}
\end{table}

% Figure block: uncomment and supply figures/overview.pdf to include an image.
% This template compiles without it.
% \begin{figure}[t]
%   \centering
%   \includegraphics[width=0.8\linewidth]{overview.pdf}
%   \caption{System overview.}
%   \label{fig:overview}
% \end{figure}
% Then cross-reference with: Figure~\ref{fig:overview}.

% =====================================================================
\section{Discussion}
\label{sec:discussion}
Interpret the results from Section~\ref{sec:results}, compare with prior work
\citep{smith2020example, brown2021inproc}, and state limitations and threats to
validity.

% =====================================================================
\section{Conclusion}
\label{sec:conclusion}
Restate the contribution and outline future work.

% --- Bibliography (natbib) ---
\bibliographystyle{plainnat}
\bibliography{references}

\end{document}
